\documentclass[../DoAn.tex]{subfiles}
\begin{document}

Chương này tổng kết những kết quả chính mà đồ án đã đạt được, đồng thời nêu ra các hướng phát triển tiếp theo để hệ thống có thể tiến gần hơn tới khả năng triển khai trong thực tế. Nếu Chương 5 tập trung phân tích các đóng góp kỹ thuật dưới góc nhìn thiết kế, thì Chương 6 đặt trọng tâm vào việc nhìn lại toàn bộ quá trình thực hiện dưới góc nhìn kết quả nghiên cứu và triển vọng ứng dụng.

\section{Kết luận}

Đồ án đã giải quyết bài toán xây dựng một trợ lý giáo dục cho môn Tin học THPT dựa trên dữ liệu sách giáo khoa chính thống. Hướng tiếp cận được lựa chọn là kết hợp truy xuất tri thức từ SGK với khả năng sinh nội dung của mô hình ngôn ngữ lớn, nhằm bảo đảm đầu ra vừa bám sát nguồn học liệu vừa đủ linh hoạt để phục vụ nhiều tác vụ khác nhau. Trên cơ sở đó, hệ thống đã được tổ chức thành một kiến trúc gồm tầng điều phối yêu cầu, tầng tri thức và tầng sinh nội dung, trong đó các tác vụ như hỏi đáp, sinh slide và soạn giáo án được xử lý trên cùng một nền dữ liệu thống nhất.

Về mặt phương pháp, đồ án đã xây dựng được một pipeline tương đối hoàn chỉnh, bắt đầu từ bước chuẩn hóa dữ liệu SGK, chunking theo cấu trúc bài học, biểu diễn nội dung bằng embedding tiếng Việt, lập chỉ mục trên vector database, kết hợp tìm kiếm từ khóa với tìm kiếm ngữ nghĩa, rồi đưa ngữ cảnh truy xuất được vào bước sinh câu trả lời hoặc sinh học liệu. Bên cạnh đó, hệ thống còn tích hợp các cơ chế điều phối như phân tích ngữ cảnh hội thoại, định tuyến ý định, quản lý trạng thái phiên và pipeline đa tác tử cho các tác vụ sinh nội dung dài. Những thành phần này giúp hệ thống không chỉ hoạt động như một chatbot hỏi đáp đơn thuần mà còn có khả năng xử lý các yêu cầu học tập phức hợp hơn.

Về mặt thực nghiệm, đồ án đã tách riêng việc đánh giá retrieve node và generate node để làm rõ đóng góp của từng thành phần trong pipeline. Với benchmark retrieval gồm 600 câu hỏi bao phủ 185 bài học, cấu hình tốt nhất của hệ thống đạt Hit@1 bằng 0,8183, Hit@10 bằng 0,9833 và MRR@10 bằng 0,8856. Kết quả này cho thấy cơ chế truy xuất được đề xuất có khả năng đưa đúng bài học mục tiêu vào các vị trí đầu của danh sách kết quả. Đối với generate node, đánh giá trên 246 mẫu bằng RAGAS cho thấy hệ thống đạt Faithfulness 0,9757 và Context Recall 0,9593, phản ánh rằng câu trả lời sinh ra nhìn chung bám sát ngữ cảnh và ít sinh thông tin ngoài tài liệu. Các kết quả đó là cơ sở thực nghiệm để khẳng định hướng tiếp cận của đồ án là khả thi đối với bài toán xây dựng trợ lý học tập bám sát SGK Tin học THPT.

Tuy vậy, đồ án vẫn còn một số giới hạn. Thứ nhất, hệ thống hiện mới được phát triển và đánh giá chủ yếu như một nguyên mẫu nghiên cứu, chưa được hoàn thiện thành một ứng dụng web có khả năng vận hành ổn định cho số lượng lớn người dùng thực. Thứ hai, miền dữ liệu mới giới hạn trong môn Tin học THPT, nên khả năng khái quát sang các môn học khác mới dừng ở mức tiềm năng kiến trúc. Thứ ba, một số thành phần vẫn phụ thuộc vào dịch vụ mô hình bên ngoài, do đó cần tiếp tục được tối ưu thêm nếu muốn triển khai ở quy mô rộng trong môi trường giáo dục thực tế.

\section{Hướng phát triển}

Hướng phát triển đầu tiên là hoàn thiện hệ thống thành một ứng dụng web đầy đủ thay vì dừng ở mức nguyên mẫu nghiên cứu. Ở hướng này, kiến trúc hiện tại cần được mở rộng thêm một lớp ứng dụng dành cho người dùng cuối, với cơ chế xác thực, phân quyền, lưu lịch sử tương tác và quản lý tài khoản một cách nhất quán. Bên cạnh kho tri thức dùng cho retrieval, hệ thống cũng cần một cơ sở dữ liệu nghiệp vụ riêng để quản lý người dùng, phiên học tập, tiến độ làm quiz, lịch sử tạo slide hoặc giáo án, cũng như các thông tin cấu hình cá nhân. Việc tách rõ cơ sở dữ liệu phục vụ tri thức và cơ sở dữ liệu phục vụ nghiệp vụ sẽ giúp hệ thống dễ bảo trì hơn và phù hợp hơn với yêu cầu triển khai thực tế.

Hướng phát triển thứ hai là thiết kế lại hệ thống theo hướng hỗ trợ nhiều người dùng đồng thời. Trong môi trường trường học thực, một hệ thống như vậy không chỉ phục vụ một người dùng đơn lẻ mà có thể phải xử lý nhiều yêu cầu song song từ học sinh và giáo viên ở các lớp khác nhau. Điều này đòi hỏi bổ sung các cơ chế như quản lý session độc lập theo người dùng, điều phối hàng đợi tác vụ dài, tối ưu tài nguyên cho retrieval và reranking, cũng như theo dõi hiệu năng ở mức hệ thống. Khi đó, các thành phần hiện tại như session state, trace và workflow đa tác tử có thể trở thành nền tảng để mở rộng theo hướng đa người dùng và triển khai theo chiều ngang.

Hướng phát triển thứ ba là đưa hệ thống vào bối cảnh sử dụng thực tế tại các trường THPT ở Việt Nam. Muốn làm được điều này, ngoài việc hoàn thiện hạ tầng phần mềm, hệ thống cần được kiểm chứng thêm thông qua các đợt thử nghiệm với học sinh và giáo viên thật. Các thử nghiệm như vậy sẽ cho phép đánh giá không chỉ độ chính xác kỹ thuật của retrieval hay generation, mà còn cả mức độ hữu ích sư phạm, mức độ dễ sử dụng và tác động thực tế tới quá trình dạy và học. Đây là bước cần thiết để chuyển đồ án từ một kết quả nghiên cứu khả thi sang một giải pháp có thể ứng dụng trong môi trường giáo dục.

Hướng phát triển thứ tư là mở rộng miền tri thức sang nhiều môn học khác ngoài Tin học. Về nguyên tắc, kiến trúc mà đồ án đề xuất không phụ thuộc tuyệt đối vào nội dung môn Tin học, mà phụ thuộc chủ yếu vào khả năng xây dựng được một kho tri thức có cấu trúc đủ tốt từ sách giáo khoa của từng môn. Do đó, nếu bổ sung được dữ liệu của các môn như Toán, Vật lí hoặc Hóa học và điều chỉnh lại một số chiến lược retrieval cũng như prompt sinh nội dung, hệ thống có thể mở rộng thành một trợ lý học tập liên môn. Khi đó, giá trị của đồ án sẽ không chỉ nằm ở một ứng dụng cho môn Tin học, mà còn ở một khung kiến trúc có thể tái sử dụng cho nhiều bài toán giáo dục khác nhau.

Cuối cùng, một hướng phát triển có ý nghĩa dài hạn là tăng cường mức độ cá nhân hóa cho người học. Hệ thống hiện đã có nền tảng quản lý trạng thái phiên và lưu vết các hoạt động học tập cơ bản. Trong tương lai, nếu mở rộng thêm phần hồ sơ người học, lịch sử câu hỏi, các lỗi thường gặp và tiến độ theo từng chủ đề, hệ thống có thể tiến tới việc gợi ý nội dung ôn tập, điều chỉnh mức độ khó của câu hỏi hoặc đề xuất học liệu phù hợp với từng học sinh. Đây là hướng phát triển phù hợp với mục tiêu lâu dài của các hệ thống giáo dục thông minh, đồng thời cũng là bước tiếp theo tự nhiên của kiến trúc hiện tại.

\end{document}
